% abtex2-modelo-artigo.tex, v-1.9.2 laurocesar
% Copyright 2012-2014 by abnTeX2 group at http://abntex2.googlecode.com/ 
%

% ------------------------------------------------------------------------
% ------------------------------------------------------------------------
% abnTeX2: Modelo de Artigo Acadêmico em conformidade com
% ABNT NBR 6022:2003: Informação e documentação - Artigo em publicação 
% periódica científica impressa - Apresentação
% ------------------------------------------------------------------------
% ------------------------------------------------------------------------

\documentclass[
	% -- opções da classe memoir --
	article,			% indica que é um artigo acadêmico
	11pt,				% tamanho da fonte
	oneside,			% para impressão apenas no verso. Oposto a twoside
	a4paper,			% tamanho do papel. 
	% -- opções da classe abntex2 --
	%chapter=TITLE,		% títulos de capítulos convertidos em letras maiúsculas
	%section=TITLE,		% títulos de seções convertidos em letras maiúsculas
	%subsection=TITLE,	% títulos de subseções convertidos em letras maiúsculas
	%subsubsection=TITLE % títulos de subsubseções convertidos em letras maiúsculas
	% -- opções do pacote babel --
	english,			% idioma adicional para hifenização
	brazil,				% o último idioma é o principal do documento
	sumario=tradicional
	]{abntex2}


% ---
% PACOTES
% ---

% ---
% Pacotes fundamentais 
% ---
\usepackage{lmodern}			% Usa a fonte Latin Modern
\usepackage[T1]{fontenc}		% Selecao de codigos de fonte.
\usepackage[utf8]{inputenc}		% Codificacao do documento (conversão automática dos acentos)
\usepackage{indentfirst}		% Indenta o primeiro parágrafo de cada seção.
\usepackage{nomencl} 			% Lista de simbolos
\usepackage{color}				% Controle das cores
\usepackage{graphicx}			% Inclusão de gráficos
\usepackage{microtype} 			% para melhorias de justificação
% ---
		
% ---
% Pacotes adicionais, usados apenas no âmbito do Modelo Canônico do abnteX2
% ---
\usepackage{lipsum}				% para geração de dummy text
% ---
		
% ---
% Pacotes de citações
% ---
\usepackage[brazilian,hyperpageref]{backref}	 % Paginas com as citações na bibl
\usepackage[alf]{abntex2cite}	% Citações padrão ABNT
% ---

% ---
% Configurações do pacote backref
% Usado sem a opção hyperpageref de backref
\renewcommand{\backrefpagesname}{Citado na(s) página(s):~}
% Texto padrão antes do número das páginas
\renewcommand{\backref}{}
% Define os textos da citação
\renewcommand*{\backrefalt}[4]{
	\ifcase #1 %
		Nenhuma citação no texto.%
	\or
		Citado na página #2.%
	\else
		Citado #1 vezes nas páginas #2.%
	\fi}%
% ---

% ---
% Informações de dados para CAPA e FOLHA DE ROSTO
% ---
\titulo{Análise Comparativa de Métodos para Detecção de\\ Face Anti-Spoofing: LBP, Deep Learning e Histogramas de Cor}
\autor{Fernando Daniel Marelino \and Ícaro Travain Darwich da Rocha \and Marcos Aquino \and Mateus Vespasiano de Castro}
\local{São Paulo, Brasil}
\data{Dezembro de 2025}
% ---

% ---
% Configurações de aparência do PDF final

% alterando o aspecto da cor azul
\definecolor{blue}{RGB}{41,5,195}

% informações do PDF
\makeatletter
\hypersetup{
    pdftitle={Analise Comparativa de Metodos para Deteccao de Face Anti-Spoofing},
    pdfauthor={Fernando Daniel Marelino; Icaro Travain; Marcos Aquino; Mateus Vespasiano},
    pdfsubject={Deteccao de Face Anti-Spoofing com Machine Learning},
    pdfcreator={LaTeX with abnTeX2},
    pdfkeywords={face anti-spoofing}{deep learning}{SVM}{histogramas de cor}{segurança da informação}, 
    colorlinks=true,       		% false: boxed links; true: colored links
    linkcolor=blue,          	% color of internal links
    citecolor=blue,        		% color of links to bibliography
    filecolor=magenta,      		% color of file links
    urlcolor=blue,
    bookmarksdepth=4
}
\makeatother

% Provisório: evita "Undefined control sequence \theforeigntitle"
% quando a classe tenta montar metadados em algumas versões do abntex2.
% Se a sua versão do abntex define esses comandos, estas linhas não atrapalham.
\providecommand{\theforeigntitle}{}
\providecommand{\foreigntitle}{}
\makeatletter
% Outros fallbacks úteis para prevenir "Undefined control sequence"
\providecommand{\thetitle}{}
\providecommand{\theauthor}{}
\providecommand{\theaffiliation}{}
\providecommand{\thekeywords}{}
\providecommand{\resumoname}{}
\providecommand{\@thanks}{}
\providecommand{\@bs@thanks}{}
\makeatother
% --- 

% ---
% compila o indice
% ---
\makeindex
% ---

% ---
% Altera as margens padrões
% ---
\setlrmarginsandblock{3cm}{3cm}{*}
\setulmarginsandblock{3cm}{3cm}{*}
\checkandfixthelayout
% ---

% --- 
% Espaçamentos entre linhas e parágrafos 
% --- 

% O tamanho do parágrafo é dado por:
\setlength{\parindent}{1.3cm}

% Controle do espaçamento entre um parágrafo e outro:
\setlength{\parskip}{0.2cm}  % tente também \onelineskip

% Espaçamento simples
\SingleSpacing

% ----
% Início do documento
% ----
\begin{document}

% Retira espaço extra obsoleto entre as frases.
\frenchspacing 

% ----------------------------------------------------------
% ELEMENTOS PRÉ-TEXTUAIS
% ----------------------------------------------------------

%---
%
% Se desejar escrever o artigo em duas colunas, descomente a linha abaixo
% e a linha com o texto ``FIM DE ARTIGO EM DUAS COLUNAS''.
% \twocolumn[    		% INICIO DE ARTIGO EM DUAS COLUNAS
%
%---
% página de titulo
\maketitle

% resumo em português
\begin{resumoumacoluna}
    Este trabalho apresenta uma análise comparativa de três métodos distintos para detecção de ataques de apresentação facial (\textit{face anti-spoofing}), um problema crítico em sistemas de autenticação biométrica. Foram avaliados: (1) Local Binary Patterns (LBP) com classificador SVM, (2) Transfer Learning utilizando a arquitetura VGG16 pré-treinada no ImageNet, e (3) Histogramas de cor nos espaços HSV e YCbCr combinados com SVM. Os experimentos foram conduzidos no dataset NUAA, amplamente utilizado na literatura. ... [falar brevemente sobre os resultados e concluir ao final]

    \vspace{\onelineskip}

    \noindent
    \textbf{Palavras-chave}: face anti-spoofing. NUAA. autenticação biométrica.
\end{resumoumacoluna}

% ]  				% FIM DE ARTIGO EM DUAS COLUNAS
% ---

% ----------------------------------------------------------
% ELEMENTOS TEXTUAIS
% ----------------------------------------------------------
\textual

% ----------------------------------------------------------
% Introdução
% ----------------------------------------------------------
\section*{Introdução}
\addcontentsline{toc}{section}{Introdução}

A autenticação biométrica baseada em reconhecimento facial tem se tornado cada vez mais prevalente em sistemas de segurança, desde o desbloqueio de smartphones até controle de acesso em ambientes corporativos e governamentais. Entretanto, essa crescente adoção também ampliou a superfície de ataque para tentativas de falsificação, onde adversários tentam enganar sistemas de reconhecimento facial utilizando fotografias impressas, vídeos reproduzidos em telas ou máscaras 3D.

Os ataques de apresentação facial, também conhecidos como \textit{face spoofing attacks}, representam uma ameaça significativa à integridade de sistemas biométricos \cite{galbally2014}. A detecção desses ataques, denominada \textit{Presentation Attack Detection} (PAD) ou \textit{face anti-spoofing}, tornou-se uma área de pesquisa ativa em visão computacional e segurança da informação.

Diversas abordagens têm sido propostas na literatura para detectar ataques de apresentação facial. Métodos baseados em textura, como \textit{Local Binary Patterns} (LBP), exploram diferenças de micro-texturas entre faces reais e falsificadas \cite{chingovska2012}. Abordagens baseadas em análise de cor \cite{DBLP:journals/corr/BoulkenafetKH15} aproveitam as distorções cromáticas introduzidas pelo processo de impressão ou exibição em telas. Mais recentemente, técnicas de \textit{deep learning} baseadas em transferência de aprendizado e em estratégias de fusão/seleção de \textit{deep features} têm demonstrado resultados promissores, aumentando robustez e acurácia em tarefas de reconhecimento visual \cite{rashid2020sustainable}.

Este trabalho apresenta uma análise comparativa de três métodos distintos para detecção de face anti-spoofing: (1) Local Binary Patterns combinado com Support Vector Machine (LBP + SVM), representando a abordagem clássica baseada em textura; (2) Transfer Learning com VGG16, representando o paradigma de \textit{deep learning}; e (3) Histogramas de cor nos espaços HSV e YCbCr combinados com SVM, explorando características cromáticas. Os experimentos foram conduzidos no dataset NUAA \cite{nuaa2010,tan2010}, um \textit{benchmark} amplamente utilizado para avaliação de métodos de anti-spoofing.

O objetivo principal é avaliar comparativamente o desempenho dessas três abordagens em termos de acurácia, \textit{Half Total Error Rate} (HTER) e \textit{Equal Error Rate} (EER), identificando vantagens e limitações de cada método. Os resultados obtidos contribuem para o entendimento de quais características são mais efetivas para a detecção de ataques de impressão fotográfica.

% ----------------------------------------------------------
% Seção de explicações
% ----------------------------------------------------------
\section{Revisão da Literatura}


\section{Metodologia}

% ---
% Finaliza a parte no bookmark do PDF, para que se inicie o bookmark na raiz
% ---
\bookmarksetup{startatroot}% 
% ---

% ---
% Seção de Experimentos
% ---
\section{Experimentos e Resultados}

% ---
% Conclusão
% ---
\section*{Conclusão}
\addcontentsline{toc}{section}{Conclusão}

% ----------------------------------------------------------
% ELEMENTOS PÓS-TEXTUAIS
% ----------------------------------------------------------
\postextual

% ---
% Título e resumo em língua estrangeira
% ---

% \twocolumn[    		% INICIO DE ARTIGO EM DUAS COLUNAS

% titulo em inglês
\titulo{Comparative Analysis of Face Anti-Spoofing Detection Methods: LBP, Deep Learning and Color Histograms}
\emptythanks
\maketitle

% resumo em português
\renewcommand{\resumoname}{Abstract}
\begin{resumoumacoluna}
 \begin{otherlanguage*}{english}
   This work presents a comparative analysis of three distinct methods for face anti-spoofing detection, a critical problem in biometric authentication systems. The following methods were evaluated: (1) Local Binary Patterns (LBP) with SVM classifier, (2) Transfer Learning using VGG16 architecture pre-trained on ImageNet, and (3) Color histograms in HSV and YCbCr color spaces combined with SVM. The experiments were conducted on the NUAA dataset, widely used in the literature. Method 3 (Color Histograms + SVM) achieved the best results, reaching 99.13\% accuracy, 0.98\% HTER, and 0.66\% EER. The results demonstrate that color-based features are highly discriminative for distinguishing real faces from photo print attacks, outperforming texture-based and deep learning approaches in this specific scenario.

   \vspace{\onelineskip}
 
   \noindent
   \textbf{Keywords}: face anti-spoofing. presentation attack detection. machine learning. SVM. deep learning. biometrics.
 \end{otherlanguage*}  
\end{resumoumacoluna}

% ]  				% FIM DE ARTIGO EM DUAS COLUNAS
% ---

% ----------------------------------------------------------
% Referências bibliográficas
% ----------------------------------------------------------
\bibliographystyle{abntex2-alf}
\bibliography{abntex2-modelo-references}

\end{document}
